
%package list
\documentclass{article}
\usepackage{amsmath}
\usepackage[top=3cm, bottom=3cm, outer=3cm, inner=3cm]{geometry}
\usepackage{multicol}
\usepackage{graphicx}
\usepackage{url}
%\usepackage{cite}
\usepackage{hyperref}
\usepackage{array}
%\usepackage{multicol}
\newcolumntype{x}[1]{>{\centering\arraybackslash\hspace{0pt}}p{#1}}
\usepackage{natbib}
\usepackage{pdfpages}
\usepackage{multirow}
\usepackage[normalem]{ulem}
\useunder{\uline}{\ul}{}
\usepackage{svg}
\usepackage{xcolor}
\usepackage{minted}

%\usepackage{booktabs}
\usepackage{caption}
\usepackage{subcaption}
\usepackage{float}
\usepackage{array}

\newcolumntype{M}[1]{>{\centering\arraybackslash}m{#1}}
\newcolumntype{N}{@{}m{0pt}@{}}

%%%%%%%%%%%%%%%%%%%%%%%%%%%%%%%%%%%%%%%%%%%%%%%%%%%%%%%%%%%%%%%%%%%%%%%%%%%%
%%%%%%%%%%%%%%%%%%%%%%%%%%%%%%%%%%%%%%%%%%%%%%%%%%%%%%%%%%%%%%%%%%%%%%%%%%%%
\newcommand{\itemEmail}{ccapatinta@unsa.edu.pe}
\newcommand{\itemStudent}{Capatinta Almanza Camilo B. A.}
\newcommand{\itemCourse}{Lab. Análisis y Diseño de Algoritmos}
\newcommand{\itemCourseCode}{1702231}
\newcommand{\itemSemester}{IV}
\newcommand{\itemUniversity}{Universidad Nacional de San Agustín de Arequipa}
\newcommand{\itemFaculty}{Facultad de Ingeniería de Producción y Servicios}
\newcommand{\itemDepartment}{Departamento Académico de Ingeniería de Sistemas e Informática}
\newcommand{\itemSchool}{Escuela Profesional de Ingeniería de Sistemas}
\newcommand{\itemAcademic}{2026 - B}
\newcommand{\itemInput}{24/09/2026}
\newcommand{\itemOutput}{28/09/2026}
\newcommand{\itemPracticeNumber}{04}
\newcommand{\itemTheme}{Teoría de Grafos. Problema del Camino Más Corto: Dijkstra y Floyd-Warshall}
%%%%%%%%%%%%%%%%%%%%%%%%%%%%%%%%%%%%%%%%%%%%%%%%%%%%%%%%%%%%%%%%%%%%%%%%%%%%
%%%%%%%%%%%%%%%%%%%%%%%%%%%%%%%%%%%%%%%%%%%%%%%%%%%%%%%%%%%%%%%%%%%%%%%%%%%%

\usepackage[english,spanish]{babel}
\usepackage[utf8]{inputenc}
\AtBeginDocument{\selectlanguage{spanish}}
\renewcommand{\figurename}{Figura}
\renewcommand{\refname}{Referencias}
\renewcommand{\tablename}{Tabla} %esto no funciona cuando se usa babel
\AtBeginDocument{%
	\renewcommand\tablename{Tabla}
}

\usepackage{fancyhdr}
\pagestyle{fancy}
\fancyhf{}
\setlength{\headheight}{40.51407pt}
\renewcommand{\headrulewidth}{1pt}
\renewcommand{\footrulewidth}{1pt}
\fancyhead[L]{\raisebox{-0.2\height}{\includegraphics[width=3cm]{../img/logo_episunsa.png}}}
\fancyhead[C]{\fontsize{7}{7}\selectfont	\itemUniversity \\ \itemFaculty \\ \itemDepartment \\ \itemSchool \\ \textbf{\itemCourse}}
\fancyhead[R]{\raisebox{-0.2\height}{\includegraphics[width=1.2cm]{../img/logo_abet}}}
\fancyfoot[L]{Estudiante Capatinta C.}
\fancyfoot[C]{\itemCourse}
\fancyfoot[R]{Página \thepage}

% para el codigo fuente

\usepackage{color, colortbl}
\definecolor{dkgreen}{rgb}{0,0.6,0}
\definecolor{gray}{rgb}{0.5,0.5,0.5}
\definecolor{mauve}{rgb}{0.58,0,0.82}
\definecolor{codebackground}{rgb}{0.95, 0.95, 0.92}
\setminted{%
  fontsize=\small,
  breaklines,
  breakanywhere,
  linenos,
  numbersep=6pt,
  frame=lines,
  framesep=3mm,
  baselinestretch=1,
  bgcolor=codebackground,
  tabsize=3
}
\definecolor{tablebackground}{rgb}{0.8, 0, 0}

\begin{document}
	
	\vspace*{10px}
	
	\begin{center}
		\parbox{0.8\textwidth}{\centering\fontsize{17}{20}\selectfont\textbf{Informe de laboratorio \itemPracticeNumber}}
	\end{center}
	\begin{center}
		\parbox{0.8\textwidth}{\centering\textbf{\Large Tema: \itemTheme}}
	\end{center}
	%\vspace*{0.5cm}	

	\begin{flushright}
		\begin{tabular}{|M{2.5cm}|N|}
			\hline 
			\rowcolor{tablebackground}
			\color{white} \textbf{Nota}  \\
			\hline 
			     \\[30pt]
			\hline 			
		\end{tabular}
	\end{flushright}	

	\begin{table}[h!]
		\begin{tabular}{|x{4.7cm}|x{4.8cm}|x{4.8cm}|}        
			\hline 
			\rowcolor{tablebackground}
			\color{white} \textbf{Estudiante} & \color{white}\textbf{Escuela}  & \color{white}\textbf{Asignatura}   \\
			\hline 
			{\itemStudent \par \itemEmail} & \itemSchool & {\itemCourse \par Semestre: \itemSemester \par Código: \itemCourseCode}     \\
			\hline 			
		\end{tabular}
	\end{table}		
	
	\begin{table}[H]
		\begin{tabular}{|x{4.7cm}|x{4.8cm}|x{4.8cm}|}
			\hline 
			\rowcolor{tablebackground}
			\color{white}\textbf{Laboratorio} & \color{white}\textbf{Tema}  & \color{white}\textbf{Duración}   \\
			\hline 
			\itemPracticeNumber & \itemTheme & 02 horas   \\
			\hline 
		\end{tabular}
	\end{table}
	
	\begin{table}[H]
		\begin{tabular}{|x{4.7cm}|x{4.8cm}|x{4.8cm}|}
			\hline 
			\rowcolor{tablebackground}
			\color{white}\textbf{Semestre académico} & \color{white}\textbf{Fecha de inicio}  & \color{white}\textbf{Fecha de entrega}   \\
			\hline 
			\itemAcademic & \itemInput &  \itemOutput  \\
			\hline 
		\end{tabular}
	\end{table}
	
\section{Ejercicios/Problemas propuestos}

\subsection{Ejercicio 1}

/relleno

\inputminted{cpp}{../e-n/e1.cpp}

\begin{itemize}
    \item /relleno
\end{itemize}

\IfFileExists{../img/e1.png}{
    \begin{figure}[H]
        \centering
        \includegraphics[width=0.8\textwidth]{../img/e1.png}
        \caption{Ejecución de e1.cpp}
    \end{figure}
}{
    \begin{center}
        \textit{Imagen de ejecución pendiente: ../img/e1.png}
    \end{center}
}

\subsection{Ejercicio 2}

/relleno

\inputminted{cpp}{../e-n/e2.cpp}

\begin{itemize}
    \item /relleno
\end{itemize}

\IfFileExists{../img/e2.png}{
    \begin{figure}[H]
        \centering
        \includegraphics[width=0.9\textwidth]{../img/e2.png}
        \caption{Ejecución de e2.cpp}
    \end{figure}
}{
    \begin{center}
        \textit{Imagen de ejecución pendiente: ../img/e2.png}
    \end{center}
}

\subsection{Ejercicio 3}

/relleno

\inputminted{cpp}{../e-n/e3.cpp}

\begin{itemize}
    \item /relleno
\end{itemize}

\IfFileExists{../img/e3.png}{
    \begin{figure}[H]
        \centering
        \includegraphics[width=0.8\textwidth]{../img/e3.png}
        \caption{Ejecución de e3.cpp}
    \end{figure}
}{
    \begin{center}
        \textit{Imagen de ejecución pendiente: ../img/e3.png}
    \end{center}
}

\subsection{Ejercicio 4}

/relleno

\inputminted{cpp}{../e-n/e4.cpp}

\begin{itemize}
    \item /relleno
\end{itemize}

\IfFileExists{../img/e4.png}{
    \begin{figure}[H]
        \centering
        \includegraphics[width=0.8\textwidth]{../img/e4.png}
        \caption{Ejecución de e4.cpp}
    \end{figure}
}{
    \begin{center}
        \textit{Imagen de ejecución pendiente: ../img/e4.png}
    \end{center}
}

\subsection{Ejercicio 5}

/relleno

\inputminted{cpp}{../e-n/e5.cpp}

\begin{itemize}
    \item /relleno
\end{itemize}

\IfFileExists{../img/e5.png}{
    \begin{figure}[H]
        \centering
        \includegraphics[width=0.8\textwidth]{../img/e5.png}
        \caption{Ejecución de e5.cpp}
    \end{figure}
}{
    \begin{center}
        \textit{Imagen de ejecución pendiente: ../img/e5.png}
    \end{center}
}

\section{Cuestionario}

\begin{enumerate}
    \item ¿Por qué el algoritmo de Dijkstra no funciona correctamente si el grafo tiene aristas con pesos negativos?

    \item ¿Qué ventaja tiene Floyd-Warshall frente a ejecutar Dijkstra una vez por cada vértice como origen?

    \item ¿Cómo cambia la complejidad de Dijkstra al usar una cola de prioridad en lugar de búsqueda lineal del mínimo?

    \item ¿Qué estructura de datos usarías para reconstruir la ruta completa, no solo su distancia total?

    \item ¿En qué situaciones reales elegirías Floyd-Warshall en lugar de Dijkstra?
\end{enumerate}

\section{Conclusiones}
% Pendiente de completar.

	%\begin{itemize}	
	%	\item Se creo el archivo \textbf{.gitignore} para no considerar los archivos %\textbf{*.class} que son innecesarios hacer seguimiento.
	%\end{itemize}
	%\begin{lstlisting}[language=bash]
	%	$ vim lab01/.gitignore
	%\end{lstlisting}
	%\begin{lstlisting}[language=bash]
	%	*.class
	%\end{lstlisting}
	%\begin{lstlisting}[language=bash]
	%	$ git add .
	%	$ git commit -m "Creando .gitignore para archivos *.class"			
	%	$ git push -u origin main
	%\end{lstlisting}
	
	%\begin{lstlisting}[language=bash]
	%	$ vim lab01/Insertion.java
	%\end{lstlisting}
	
    
%\clearpage
%\bibliographystyle{apalike}
%\bibliographystyle{IEEEtranN}
%\bibliography{bibliography}
			
\end{document}\noindent\textbf{Enunciado:} 
